\section{The principles of inclusion and exclusion}

\subsection{The size version}

\begin{theorem}[Size version of the PIE]
\label{thm.pie.1}
\lean{AlgebraicCombinatorics.InclusionExclusion.pie_size_version}
\leantarget
\leanok
Let $n \in \mathbb{N}$. Let $U$ be a finite set. Let $A_1, A_2, \ldots, A_n$ be $n$ subsets of $U$. Then,
\[
\left(\text{\# of } u \in U \text{ that satisfy } u \notin A_i \text{ for all } i \in [n]\right)
= \sum_{I \subseteq [n]} (-1)^{|I|} \left(\text{\# of } u \in U \text{ that satisfy } u \in A_i \text{ for all } i \in I\right).
\]
Equivalently,
\[
\left| U \setminus (A_1 \cup A_2 \cup \cdots \cup A_n) \right|
= \sum_{I \subseteq [n]} (-1)^{|I|} \left| \bigcap_{i \in I} A_i \right|,
\]
where $\bigcap_{i \in I} A_i$ denotes $\{u \in U \mid u \in A_i \text{ for all } i \in I\}$
(so that $\bigcap_{i \in \varnothing} A_i = U$).
\end{theorem}

\begin{proof}
This is the inclusion-exclusion cardinality formula from Mathlib.
\end{proof}

\subsection{Alternative formulations of the size version}

\begin{lemma}[Complement-of-union formulation]
\label{lem.pie.size-version-prime}
\lean{AlgebraicCombinatorics.InclusionExclusion.pie_size_version'}
\leanhelper
\leanok
The complement of the union $A_1 \cup A_2 \cup \cdots \cup A_n$ satisfies
\[
\left| (A_1 \cup A_2 \cup \cdots \cup A_n)^c \right|
= \sum_{I \subseteq [n]} (-1)^{|I|}\, \left|\bigcap_{i \in I} A_i\right|.
\]
This is the more common textbook formulation of the PIE.
\end{lemma}

\begin{proof}
Rewrite $(A_1 \cup \cdots \cup A_n)^c = \bigcap_i A_i^c$, then apply
Theorem~\ref{thm.pie.1}.
\end{proof}

\begin{lemma}[PIE specialized to $\{0, 1, \ldots, n{-}1\}$]
\label{lem.pie.size-version-fin}
\lean{AlgebraicCombinatorics.InclusionExclusion.pie_size_version_fin}
\leanhelper
\leanok
Let $\alpha$ be a finite type and $A_0, A_1, \ldots, A_{n-1}$ be $n$ subsets of $\alpha$
indexed by $\{0, 1, \ldots, n{-}1\}$. Then
\[
\#\{u \in \alpha \mid u \notin A_i\ \forall\, i \in \{0, \ldots, n{-}1\}\}
= \sum_{I \subseteq \{0, \ldots, n{-}1\}} (-1)^{|I|}\, \#\!\left(\bigcap_{i \in I} A_i\right).
\]
\end{lemma}

\begin{proof}
This is the specialization of Theorem~\ref{thm.pie.1} to the index set
$\{0, 1, \ldots, n{-}1\}$.
\end{proof}

\begin{lemma}[Rule-breaking interpretation of PIE]
\label{lem.pie.rule-breaking}
\lean{AlgebraicCombinatorics.InclusionExclusion.pie_rule_breaking}
\leanhelper
\leanok
Given $n$ ``rules'' encoded as subsets $A_0, \ldots, A_{n-1}$ of a finite type
$\alpha$ (where $A_i$ is the set of elements satisfying rule $i$), the number
of elements violating all rules is
\[
\#\{u \in \alpha \mid u \notin A_i\ \forall\, i\}
= \sum_{I \subseteq \{0, \ldots, n{-}1\}} (-1)^{|I|}\,
\#\{u \in \alpha \mid u \in A_i\ \forall\, i \in I\}.
\]
\end{lemma}

\begin{proof}
Rewrite both sides in terms of infima of complements and intersections,
then apply Lemma~\ref{lem.pie.size-version-fin}.
\end{proof}

\subsection{Counting surjections}

\begin{definition}[Number of surjections]
\label{def.pie.numSurj}
\lean{AlgebraicCombinatorics.InclusionExclusion.numSurj}
\leanhelper
\leanok
The \emph{number of surjective maps} from $[m]$ to $[n]$
is defined as
\[
\mathrm{numSurj}(m, n) := \#\{f\colon [m] \to [n] \mid f \text{ is surjective}\}.
\]
\end{definition}

\begin{theorem}
\label{thm.pie.count-sur}
\lean{AlgebraicCombinatorics.InclusionExclusion.numSurj_formula}
\leantarget
\leanok
Let $n, m \in \mathbb{N}$. Then,
\[
\left(\text{\# of surjective maps } f\colon [m] \to [n]\right)
= \sum_{k=0}^{n} (-1)^k \binom{n}{k} (n-k)^m.
\]
\end{theorem}

\begin{proof}
A map $f\colon [m] \to [n]$ is surjective if and only if it takes each $i \in [n]$
as a value. Thus, imposing $n$ rules on a map $f\colon [m] \to [n]$, where
rule $i$ says ``thou shalt not take $i$ as a value,'' the surjective maps
are precisely the maps that violate all $n$ rules. By the PIE
(Theorem~\ref{thm.pie.1}):
\begin{align*}
& \left(\text{\# of surjective maps } f\colon [m] \to [n]\right) \\
& = \sum_{I \subseteq [n]} (-1)^{|I|}
\left(\text{\# of maps } f\colon [m] \to [n] \text{ that satisfy all rules in } I\right).
\end{align*}
For each $I \subseteq [n]$, the maps satisfying all rules in $I$ are
the maps from $[m]$ to $[n] \setminus I$, of which there are
$(n - |I|)^m$.
Substituting and splitting the sum by $|I| = k$:
\[
\sum_{k=0}^{n} \binom{n}{k} (-1)^k (n-k)^m.
\]

In Lean, we relate the surjection count to the PIE via the ``avoid'' sets
$S_i = \{f \mid i \notin \operatorname{im} f\}$, apply
the PIE (Theorem~\ref{thm.pie.1}), substitute the cardinality formula
for each subset, and rewrite the sum over the powerset as a sum over cardinalities.
\end{proof}

\begin{lemma}[Cardinality of functions avoiding a set]
\label{lem.pie.card_functions_avoiding_set}
\lean{AlgebraicCombinatorics.InclusionExclusion.card_functions_avoiding_set}
\leanhelper
\leanok
Let $m, n \in \mathbb{N}$ and let $I$ be a subset of $[n]$. The number
of functions $f\colon [m] \to [n]$ such that $i \notin \operatorname{im} f$
for all $i \in I$ is $(n - |I|)^m$.
\end{lemma}

\begin{proof}
Such functions are exactly the functions from $[m]$ to $I^c$,
and $|I^c| = n - |I|$.
\end{proof}

\begin{lemma}[Cardinality of intersection of avoid-sets]
\label{lem.pie.card_inf_S_avoid}
\lean{AlgebraicCombinatorics.InclusionExclusion.card_inf_S_avoid}
\leanhelper
\leanok
For $t \subseteq [n]$, the cardinality of the intersection
$\bigcap_{i \in t} S_{\mathrm{avoid}}(i)$ is $(n - |t|)^m$,
where $S_{\mathrm{avoid}}(i) = \{f\colon [m] \to [n] \mid i \notin \operatorname{im} f\}$.
\end{lemma}

\begin{proof}
The intersection equals the set of functions whose range avoids all elements of $t$,
which by Lemma~\ref{lem.pie.card_functions_avoiding_set} has cardinality $(n - |t|)^m$.
\end{proof}

\subsection{Consequences of the surjection formula}

\begin{corollary}
\label{cor.pie.count-sur.cors}
\lean{AlgebraicCombinatorics.InclusionExclusion.surjOn_alternating_sum_eq_zero, AlgebraicCombinatorics.InclusionExclusion.surjOn_alternating_sum_eq_factorial, AlgebraicCombinatorics.InclusionExclusion.surjOn_alternating_sum_nonneg, AlgebraicCombinatorics.InclusionExclusion.surjOn_alternating_sum_dvd_factorial}
\leantarget
\leanok
Let $n \in \mathbb{N}$. Then:

\textbf{(a)} We have $\sum_{k=0}^{n} (-1)^k \binom{n}{k} (n-k)^m = 0$ for any $m \in \mathbb{N}$ satisfying $m < n$.

\textbf{(b)} We have $\sum_{k=0}^{n} (-1)^k \binom{n}{k} (n-k)^n = n!$.

\textbf{(c)} We have $\sum_{k=0}^{n} (-1)^k \binom{n}{k} (n-k)^m \geq 0$ for each $m \in \mathbb{N}$.

\textbf{(d)} We have $n! \mid \sum_{k=0}^{n} (-1)^k \binom{n}{k} (n-k)^m$ for each $m \in \mathbb{N}$.
\end{corollary}

\begin{proof}
\textbf{(a)} By Theorem~\ref{thm.pie.count-sur}, the sum equals the number of surjective
maps $f\colon [m] \to [n]$. Since $m < n$, no such surjection exists (by the
Pigeonhole Principle), so the sum is $0$.

\textbf{(b)} By Theorem~\ref{thm.pie.count-sur} with $m = n$, the sum counts surjections
$f\colon [n] \to [n]$, which are precisely the permutations of $[n]$, of which there
are $n!$.

\textbf{(c)} The sum equals the number of surjections, which is a nonnegative integer.

\textbf{(d)} The symmetric group $S_n$ acts on the set
$U := \{\text{surjective maps } f\colon [m] \to [n]\}$
by post-composition ($\sigma \cdot f = \sigma \circ f$).
This action is free: if $\sigma \circ f = f$ and $f$ is surjective, then
for each $j \in [n]$, pick $i$ with $f(i) = j$; then
$\sigma(j) = \sigma(f(i)) = f(i) = j$, so $\sigma = \mathrm{id}$
(Lemma~\ref{lem.pie.stabilizer_surj_eq_bot}).
By the orbit--stabilizer theorem, each orbit has size $n!$, so the orbits
partition $U$ into blocks of size $n!$, whence
$n! \mid |U|$ (Lemma~\ref{lem.pie.factorial_dvd_numSurj}).
\end{proof}

\subsection{Helpers for the divisibility proof}

\begin{lemma}[Post-composition preserves surjectivity]
\label{lem.pie.surjective_smul}
\lean{AlgebraicCombinatorics.InclusionExclusion.surjective_smul}
\leanhelper
\leanok
If $\sigma$ is a permutation of $[n]$ and
$f\colon [m] \to [n]$ is surjective, then
$\sigma \circ f$ is surjective.
\end{lemma}

\begin{proof}
Given $y \in [n]$, since $f$ is surjective there exists $x$ with
$f(x) = \sigma^{-1}(y)$, so $\sigma(f(x)) = y$.
\end{proof}

\begin{lemma}[Action of $S_n$ on surjections]
\label{lem.pie.permActionOnSurj}
\lean{AlgebraicCombinatorics.InclusionExclusion.permActionOnSurj}
\leanhelper
\leanok
The symmetric group $S_n$ acts on the set of surjective
functions $f\colon [m] \to [n]$ by post-composition:
$\sigma \cdot f := \sigma \circ f$.
\end{lemma}

\begin{proof}
\leanok
Well-definedness follows from Lemma~\ref{lem.pie.surjective_smul}. The action
laws $1 \cdot f = f$ and $(\sigma\tau) \cdot f = \sigma \cdot (\tau \cdot f)$ are
immediate from composition.
\end{proof}

\begin{lemma}[Free action on surjections]
\label{lem.pie.stabilizer_surj_eq_bot}
\lean{AlgebraicCombinatorics.InclusionExclusion.stabilizer_surj_eq_bot}
\leanhelper
\leanok
Let $f\colon [m] \to [n]$ be surjective. Then the
stabilizer of $f$ under the action of $S_n$ by
post-composition is trivial.
\end{lemma}

\begin{proof}
If $\sigma \circ f = f$ and $f$ is surjective, then for each $y \in [n]$,
pick $x$ with $f(x) = y$; then $\sigma(y) = \sigma(f(x)) = f(x) = y$. So $\sigma = \mathrm{id}$.
\end{proof}

\begin{lemma}[Factorial divides number of surjections]
\label{lem.pie.factorial_dvd_numSurj}
\lean{AlgebraicCombinatorics.InclusionExclusion.factorial_dvd_numSurj}
\leanhelper
\leanok
We have $n! \mid \mathrm{numSurj}(m, n)$.
\end{lemma}

\begin{proof}
By the orbit--stabilizer theorem, each orbit of the free action of $S_n$
on surjections has size $|S_n|/|\{\mathrm{id}\}| = n!$, so the orbits partition
the set of surjections into blocks of size $n!$.
\end{proof}

\subsection{Derangements}

\begin{definition}
\label{def.pie.dera}
\lean{AlgebraicCombinatorics.InclusionExclusion.Derangement}
\leantarget
\leanok
A \emph{derangement} of a set $X$ means a permutation of $X$ that has no fixed points.
\end{definition}

We write $D_n$ for the number of derangements of $[n]$.

\begin{lemma}[Equivalence to Mathlib's derangements]
\label{lem.pie.equivDerangements}
\lean{AlgebraicCombinatorics.InclusionExclusion.Derangement.equivDerangements}
\leanhelper
\leanok
The set of derangements of $\alpha$ (as defined in Definition~\ref{def.pie.dera})
is in bijection with the set of fixed-point-free permutations
from Mathlib.
\end{lemma}

\begin{proof}
\leanok
Both are defined by the same predicate
(having no fixed points) on the set of permutations of $\alpha$.
\end{proof}

\begin{lemma}[Derangements of the empty type]
\label{lem.pie.derangement-empty}
\lean{AlgebraicCombinatorics.InclusionExclusion.Derangement.id_of_isEmpty}
\leanhelper
\leanok
The identity permutation is a derangement of the empty type
(vacuously, it has no fixed points). Hence $D_0 = 1$.
\end{lemma}

\begin{proof}
\leanok
Since the type is empty, there are no elements to be fixed points.
\end{proof}

\begin{lemma}[Identity is not a derangement of nonempty types]
\label{lem.pie.derangement-nonempty}
\lean{AlgebraicCombinatorics.InclusionExclusion.Derangement.id_not_derangement}
\leanhelper
\leanok
If $\alpha$ is nonempty, then no derangement of $\alpha$ equals the identity
permutation. In particular $D_1 = 0$.
\end{lemma}

\begin{proof}
\leanok
If $d$ is a derangement with $d = \mathrm{id}$, pick any $x \in \alpha$;
then $d(x) = x$, contradicting the fixed-point-free property.
\end{proof}

\begin{lemma}[Small derangement counts]
\label{lem.pie.derangement-small-values}
\lean{AlgebraicCombinatorics.InclusionExclusion.numDerangements_two', AlgebraicCombinatorics.InclusionExclusion.numDerangements_three}
\leanhelper
\leanok
The first few values of the derangement count are:
$D_2 = 1$ (the only derangement of $\{0,1\}$ is the swap) and
$D_3 = 2$ (the two $3$-cycles).
\end{lemma}

\begin{proof}
\leanok
By computation (verified by decision procedure in Lean).
\end{proof}

\begin{lemma}[Derangement type matches Mathlib]
\label{lem.pie.card_Derangement_eq}
\lean{AlgebraicCombinatorics.InclusionExclusion.card_Derangement_eq}
\leanhelper
\leanok
The number of derangements of $[n]$ (as defined in
Definition~\ref{def.pie.dera}) agrees with Mathlib's derangement count.
\end{lemma}

\begin{proof}
\leanok
Follows from Lemma~\ref{lem.pie.equivDerangements} and
Mathlib's computation of the derangement count.
\end{proof}

\begin{theorem}
\label{thm.pie.count-der}
\lean{AlgebraicCombinatorics.InclusionExclusion.numDerangements_formula}
\leantarget
\leanok
Let $n \in \mathbb{N}$. Then, the number of derangements of $[n]$ is
\[
D_n = \sum_{k=0}^{n} (-1)^k \binom{n}{k} (n-k)!
= n! \cdot \sum_{k=0}^{n} \frac{(-1)^k}{k!}.
\]
\end{theorem}

\begin{proof}
Set $U := S_n$ (the set of all permutations of $[n]$). Impose $n$ rules
$1, 2, \ldots, n$ on a permutation $\sigma \in U$, with rule $i$ being
``thou shalt leave the element $i$ fixed'' (i.e., rule $i$ requires
$\sigma(i) = i$). Derangements are precisely the permutations violating all
$n$ rules. By the PIE (Theorem~\ref{thm.pie.1}):
\[
D_n = \sum_{I \subseteq [n]} (-1)^{|I|}
\left(\text{\# of permutations } u \in U \text{ that satisfy all rules in } I\right).
\]
For each $I \subseteq [n]$, the permutations satisfying all rules in $I$ are
the permutations of $[n]$ that leave each $i \in I$ fixed; there are
$(n - |I|)!$ such permutations (they are essentially the permutations of the
$(n - |I|)$-element set $[n] \setminus I$).
Substituting and grouping by $|I| = k$:
\[
D_n = \sum_{k=0}^{n} \binom{n}{k} (-1)^k (n-k)!
= \sum_{k=0}^{n} (-1)^k \frac{n!}{k!}
= n! \sum_{k=0}^{n} \frac{(-1)^k}{k!}.
\]

In Lean, the first equality is derived from Mathlib's derangement sum formula
using the identity $(n-k)! \cdot \binom{n}{k} = n!/k!$.
\end{proof}

\begin{theorem}[Derangement formula, rational form]
\label{thm.pie.count-der.rat}
\lean{AlgebraicCombinatorics.InclusionExclusion.numDerangements_formula_rat}
\leanhelper
\leanok
Let $n \in \mathbb{N}$. Then,
\[
D_n = n! \cdot \sum_{k=0}^{n} \frac{(-1)^k}{k!}.
\]
(This is the same formula as in Theorem~\ref{thm.pie.count-der}, but stated
over $\mathbb{Q}$ so that the division $1/k!$ is literal.)
\end{theorem}

\begin{proof}
Convert Theorem~\ref{thm.pie.count-der} from $\mathbb{Z}$ to $\mathbb{Q}$, then
use $\binom{n}{k} (n-k)! = n!/k!$.
\end{proof}

\subsection{Euler's totient formula}

\begin{theorem}
\label{thm.pie.euler-tot}
\lean{AlgebraicCombinatorics.InclusionExclusion.totient_eq_prod_one_sub_inv}
\leantarget
\leanok
Let $c$ be a positive integer with prime factorization
$p_1^{a_1} p_2^{a_2} \cdots p_n^{a_n}$, where $p_1, p_2, \ldots, p_n$ are
distinct primes and $a_1, a_2, \ldots, a_n$ are positive integers. Then,
\[
\left(\text{\# of all } u \in [c] \text{ that are coprime to } c\right)
= c \cdot \prod_{i=1}^{n} \left(1 - \frac{1}{p_i}\right)
= \prod_{i=1}^{n} \left(p_i^{a_i} - p_i^{a_i - 1}\right).
\]
\end{theorem}

\begin{proof}
A number $u \in [c]$ is coprime to $c$ if and only if it is not divisible by any of
the prime factors $p_1, \ldots, p_n$ of $c$. Thus $u$ must break all $n$ rules
$1, 2, \ldots, n$, where rule $i$ says ``thou shalt be divisible by $p_i$.''
By the PIE (Theorem~\ref{thm.pie.1}):
\begin{align*}
\phi(c) &= \sum_{I \subseteq [n]} (-1)^{|I|}
\underbrace{\left(\text{\# of all } u \in [c] \text{ that are divisible by all }
p_i \text{ with } i \in I\right)}_{= c / \prod_{i \in I} p_i} \\
&= c \cdot \underbrace{\sum_{I \subseteq [n]} (-1)^{|I|}
\prod_{i \in I} \frac{1}{p_i}}_{= \prod_{i=1}^{n} (1 - 1/p_i)} \\
&= \underbrace{c}_{= \prod_{i=1}^{n} p_i^{a_i}} \cdot \prod_{i=1}^{n}
\left(1 - \frac{1}{p_i}\right)
= \prod_{i=1}^{n} \left(p_i^{a_i} - p_i^{a_i - 1}\right).
\end{align*}
(The factorization of the sum over subsets into a product uses the identity
$\sum_{I \subseteq [n]} \prod_{i \in I} x_i = \prod_{i=1}^{n}(1 + x_i)$.)

In Lean, this is derived from Mathlib's totient formula.
\end{proof}

\subsection{The weighted version}

\begin{theorem}[Weighted version of the PIE]
\label{thm.pie.2}
\lean{AlgebraicCombinatorics.InclusionExclusion.weighted_pie}
\leantarget
\leanok
Let $n \in \mathbb{N}$, and let $U$ be a finite set. Let $A_1, A_2, \ldots, A_n$
be $n$ subsets of $U$. Let $G$ be any additive abelian group. Let $w\colon U \to G$
be any map. Then,
\[
\sum_{\substack{u \in U \\ u \notin A_i \text{ for all } i \in [n]}} w(u)
= \sum_{I \subseteq [n]} (-1)^{|I|} \sum_{\substack{u \in U \\ u \in A_i \text{ for all } i \in I}} w(u).
\]
\end{theorem}

\begin{proof}
This is the weighted inclusion-exclusion formula from Mathlib.
\end{proof}

\begin{theorem}[Weighted PIE, explicit index set]
\label{thm.pie.2.explicit}
\lean{AlgebraicCombinatorics.InclusionExclusion.weighted_pie'}
\leanhelper
\leanok
Let $U$ be a finite set, $s$ a finite index set, $A_i$ subsets of $U$ for $i \in s$,
$G$ an additive abelian group, and $w\colon U \to G$ a weight function. Then,
\[
\sum_{\substack{u \in U \\ u \notin A_i\ \forall\, i \in s}} w(u)
= \sum_{I \subseteq s} (-1)^{|I|} \sum_{\substack{u \in U \\ u \in A_i\ \forall\, i \in I}} w(u).
\]
\end{theorem}

\begin{proof}
Specialization of Theorem~\ref{thm.pie.2} to an explicit index set.
\end{proof}

\begin{lemma}[Size PIE from the weighted PIE]
\label{lem.pie.size-from-weighted}
\lean{AlgebraicCombinatorics.InclusionExclusion.size_pie_from_weighted}
\leanhelper
\leanok
The size version of the PIE (Theorem~\ref{thm.pie.1}) follows from the weighted version
(Theorem~\ref{thm.pie.2}) by taking the weight function $w(u) = 1$ for all $u$.
Explicitly, for a finite type $U$, finite index set $s$, and subsets $A_i$ of $U$:
\[
\left|\bigcap_{i \in s} A_i^c\right|
= \sum_{I \in s.\mathrm{powerset}} (-1)^{|I|}\, \left|\bigcap_{i \in I} A_i\right|.
\]
\end{lemma}

\begin{proof}
Specialize Theorem~\ref{thm.pie.2.explicit} to $G = \mathbb{Z}$ and $w \equiv 1$.
Each weighted sum $\sum_{u \in S} w(u) = \sum_{u \in S} 1 = |S|$.
\end{proof}
